\begin{table}[!ht]
\centering
\caption{Decomposição interseccional (raça $\times$ gênero) do gap de rendimento
\emph{vs.}\ o Homem Branco (grupo de referência). Dotações: parcela explicada por
características observáveis; Retornos: parcela não explicada (discriminação). Penalidade
extra: o quanto o gap da Mulher Negra excede a soma das penalidades de raça e de gênero
isoladas (efeito interseccional puro, Crenshaw, 1989). PNAD Contínua, população completa
($N$ do grupo de referência --- homens brancos --- na última linha da nota). Entre
parênteses: erro-padrão por bootstrap em blocos por UPA.}
\label{tab:interseccional}
\resizebox{\textwidth}{!}{%
\begin{tabular}{lrrrrr}
\toprule
Grupo & Gap vs HB (\%) & Dotações (\%) & Retornos (\%) & Penal.\ extra (p.p.) & $N$ do grupo \\
\midrule
Mulher Branca & 19,7 (0,2) & −51,9 (1,0) & 151,9 (1,0) & --- & 1.423.976 \\
Homem Negro & 56,3 (0,6) & 73,4 (0,5) & 26,6 (0,5) & --- & 2.664.706 \\
Mulher Negra & 80,6 (0,6) & 34,9 (0,4) & 65,1 (0,4) & 4,6 & 1.769.917 \\
\bottomrule
\end{tabular}
}
\par\smallskip
\footnotesize\emph{Como ler:} ``Gap vs HB'' é a desvantagem de renda de cada grupo
\emph{vs.}\ o Homem Branco; Dotações $+$ Retornos $=100\%$ do gap. A Mulher Negra tem o
maior gap e ainda uma penalidade \emph{extra} que não se reduz à soma de ``ser negro'' e
``ser mulher'' --- a marca da interseccionalidade. O sinal negativo das dotações da Mulher
Branca indica que suas características observáveis \emph{superam} as do homem branco: todo
o gap vem de retornos diferenciais. Grupo de referência: 1.835.496 homens brancos.
\end{table}
